% GENERATED FILE. Edit canonical lessons, then run npm run generate:books.
% canonical-source-hash: fnv1a64:83e02ed65a439746
% canonical-lessons: ZH-C465-tongkuai, ZH-C465-saoxing, ZH-W465-xing-mood, ZH-C465-leguan, ZH-C465-beiguan, ZH-C465-zixin

\chapter{Spoiling the Fun}
\label{ch:zh-spoiling-the-fun}
\hlchaptermodality{465}

\begin{chapteropening}
\textbf{By the end of this chapter:} \emph{I can say we had a great time, that something spoiled the fun, and talk about being optimistic, gloomy or sure of myself.}

Complete the last lesson of chapter 465: I can say we had a great time, that something spoiled the fun, and talk about being optimistic, gloomy or sure of myself.
\end{chapteropening}

\section[tòngkuài]{tòngkuài — delighted; to one's heart's content}

\label{lesson:ZH-C465-tongkuai}



\begin{quote}
Before the new one: say the Mandarin for to take it out on someone, then the Mandarin for furious.
\end{quote}

\subsection*{You'll want to know: tòngkuài}
\emph{tòngkuài} — \textquotedblleft{}delighted; to one's heart's content\textquotedblright{}.

\begin{quote}
\emph{wǒmen wán le yì tiān wán de hěn tòngkuài} — We spent the whole day having fun and enjoyed every minute of it.
\end{quote}

Its characters come later; for now it is a word for the ear.

\begin{cousinweb}
\emph{tòngkuài} is built from \emph{tòng} 'to hurt' and \zh{快} 'glad', both inherited from Middle Chinese, \emph{thuwngH} and \emph{khwaejH} (Guangyun, in Baxter's transcription), which glosses the first 'illness' and 'to hurt' and the second as what pleases the heart. The sense 'thoroughly' for the first is not given there.
\end{cousinweb}

\begin{grammarlens}[title={What you've built}]
One more word for saying how it was.
\end{grammarlens}

\subsection*{Your turn}
\begin{itemize}
\raggedright
  \item \emph{Say it:} \emph{tòngkuài}
  \item \emph{Say it:} \emph{tòngkuài}, once more
  \item \emph{Say it:} \emph{wǒmen wán le yì tiān wán de hěn tòngkuài} — We spent the whole day having fun and enjoyed every minute of it.
  \item \emph{Recall:} say the Mandarin for to take it out on someone, then the Mandarin for furious, then say \emph{tòngkuài} again
\end{itemize}

\begin{tcolorbox}[breakable,colback=teal!4,colframe=teal!35!black,title={Before you move on}]
What does tòngkuài mean? (\textquotedblleft{}Delighted; to one's heart's content\textquotedblright{}.) Say it once more.
\end{tcolorbox}
\section[\texorpdfstring{\zh{扫兴} (sǎoxìng)}{扫兴 (sǎoxìng)}]{\zh{扫兴} (sǎoxìng) — deflated; with the fun spoiled}

\label{lesson:ZH-C465-saoxing}



\begin{quote}
Before the new one: say the Mandarin for furious, then the Mandarin for to one's heart's content.
\end{quote}

\subsection*{You'll want to know: \zh{扫兴}}
\textbf{\zh{扫兴}} — \emph{sǎoxìng} — \textquotedblleft{}deflated; with the fun spoiled\textquotedblright{}.

\begin{quote}
\textbf{\zh{下雨了真扫兴}} \emph{xiàyǔ le zhēn sǎoxìng} — It's raining. What a let-down.
\end{quote}

The new character, in pieces: \zh{兴} — six strokes in three pieces: at the top a dot, a second dot and a short stroke falling to the left; a long bar across; and under it \zh{八} \emph{bā}, two strokes. The next lesson puts it in your hand.

\begin{cousinweb}
\textbf{\zh{扫兴}} (\emph{sǎoxìng}) is built from \zh{扫} 'to sweep' and \zh{兴} 'high spirits', both inherited from Middle Chinese, \emph{sawX} and \emph{xingH} (Guangyun, in Baxter's transcription). The dictionary glosses the first 'to sweep away'; it gives the second reading no gloss of its own, so 'high spirits' is the usual explanation, not a settled one.
\end{cousinweb}

\begin{grammarlens}[title={What you've built}]
One more word for saying how it was.
\end{grammarlens}

\subsection*{Your turn}
\begin{itemize}
\raggedright
  \item \emph{Say it:} \emph{sǎoxìng}
  \item \emph{Say it:} \emph{sǎoxìng}, once more
  \item \emph{Say it:} \emph{xiàyǔ le zhēn sǎoxìng} — It's raining. What a let-down.
  \item \emph{Recall:} say the Mandarin for furious, then the Mandarin for to one's heart's content, then say \emph{sǎoxìng} again
\end{itemize}

\begin{tcolorbox}[breakable,colback=teal!4,colframe=teal!35!black,title={Before you move on}]
What does \zh{扫兴} mean? (\textquotedblleft{}Deflated; with the fun spoiled\textquotedblright{}.) Say it once more.
\end{tcolorbox}
\section[\texorpdfstring{\zh{兴} (xìng)}{兴 (xìng)}]{\zh{兴} — three short strokes over a bar, and \zh{八} under it}

\label{lesson:ZH-W465-xing-mood}



\begin{quote}
The word first, then the character inside it.

\begin{itemize}
\raggedright
  \item \emph{Recall:} say \emph{sǎoxìng} — \textbf{R1}, one lesson back, before it decays
\end{itemize}

\begin{quote}
\zh{八} \zh{高}
\end{quote}

\emph{Read it:} each one aloud, with nothing said first — \emph{bā}, \emph{gāo}
\end{quote}

\subsection*{Script — three short strokes over a bar, and \zh{八} under it}
\hlblockfigure{\detokenize{figures/ZH-W465-xing-mood-filmstrip.pdf}}{How \zh{兴} is written, stroke by stroke}

\begin{quote}
\zh{兴}
\end{quote}

\textbf{Three pieces, top to bottom.} Count the pieces, not the strokes.

\textbf{The top, three strokes.} A dot on the left, a second dot in the middle, both falling to the right, then a short stroke on the right falling to the left.

\textbf{The bar, one stroke.} A long bar across, from left to right, wider than everything above it.

\textbf{The bottom, \zh{八}, two strokes.} A stroke falling to the left, then a short stroke falling to the right.

3 + 1 + 2 = six.

\subsection*{Your turn}
\begin{itemize}
\raggedright
  \item \emph{Write it:} \zh{兴}, all six, with the model in front of you
  \item \emph{Cover:} write it again from memory, and count the pieces, not the strokes
\end{itemize}

\begin{tcolorbox}[breakable,colback=teal!4,colframe=teal!35!black,title={Before you move on}]
How many pieces does \zh{兴} have? (\textbf{Three: three strokes, one and two}.) What is under the bar? (\textbf{\zh{八}}.)

Source: \href{https://raw.githubusercontent.com/chanind/hanzi-writer-data/68d10a4b21150cae5e1ebbd223eed289cf32d90c/data/\%E5\%85\%B4.json}{Hanzi Writer Data \zh{兴}, pinned stroke paths and medians}.
\end{tcolorbox}
\section[lèguān]{lèguān — optimistic}

\label{lesson:ZH-C465-leguan}



\begin{quote}
Before the new one: say the Mandarin for to one's heart's content, then the Mandarin for the fun spoiled.
\end{quote}

\subsection*{You'll want to know: lèguān}
\emph{lèguān} — \textquotedblleft{}optimistic\textquotedblright{}.

\begin{quote}
\emph{wǒ nǎinai shì wǒmen jiā zuì lèguān de rén} — My grandmother is the most optimistic person in our family.
\end{quote}

Its characters come later; for now it is a word for the ear.

\begin{cousinweb}
\emph{lèguān} is built from \zh{乐} 'joy' and \emph{guān} 'to view', both inherited from Middle Chinese, \emph{lak} and \emph{kwan} (Guangyun, in Baxter's transcription), which glosses the first 'joy' and the second 'to look': looking at things with joy.
\end{cousinweb}

\begin{grammarlens}[title={What you've built}]
One more word for saying how it was.
\end{grammarlens}

\subsection*{Your turn}
\begin{itemize}
\raggedright
  \item \emph{Say it:} \emph{lèguān}
  \item \emph{Say it:} \emph{lèguān}, once more
  \item \emph{Say it:} \emph{wǒ nǎinai shì wǒmen jiā zuì lèguān de rén} — My grandmother is the most optimistic person in our family.
  \item \emph{Recall:} say the Mandarin for to one's heart's content, then the Mandarin for the fun spoiled, then say \emph{lèguān} again
\end{itemize}

\begin{tcolorbox}[breakable,colback=teal!4,colframe=teal!35!black,title={Before you move on}]
What does lèguān mean? (\textquotedblleft{}Optimistic\textquotedblright{}.) Say it once more.
\end{tcolorbox}
\section[bēiguān]{bēiguān — pessimistic}

\label{lesson:ZH-C465-beiguan}



\begin{quote}
Before the new one: say the Mandarin for the fun spoiled, then the Mandarin for optimistic.
\end{quote}

\subsection*{You'll want to know: bēiguān}
\emph{bēiguān} — \textquotedblleft{}pessimistic\textquotedblright{}.

\begin{quote}
\emph{nǐ bié zhème bēiguān} — Don't be so pessimistic.
\end{quote}

Its characters come later; for now it is a word for the ear.

\begin{cousinweb}
\emph{bēiguān} is built from \emph{bēi} 'grief' and \emph{guān} 'to view', both inherited from Middle Chinese, \emph{pij} and \emph{kwan} (Guangyun, in Baxter's transcription), which glosses the first 'pain' and the second 'to look'.
\end{cousinweb}

\begin{grammarlens}[title={What you've built}]
One more word for saying how it was.
\end{grammarlens}

\subsection*{Your turn}
\begin{itemize}
\raggedright
  \item \emph{Say it:} \emph{bēiguān}
  \item \emph{Say it:} \emph{bēiguān}, once more
  \item \emph{Say it:} \emph{nǐ bié zhème bēiguān} — Don't be so pessimistic.
  \item \emph{Recall:} say the Mandarin for the fun spoiled, then the Mandarin for optimistic, then say \emph{bēiguān} again
\end{itemize}

\begin{tcolorbox}[breakable,colback=teal!4,colframe=teal!35!black,title={Before you move on}]
What does bēiguān mean? (\textquotedblleft{}Pessimistic\textquotedblright{}.) Say it once more.
\end{tcolorbox}
\section[\texorpdfstring{\zh{自信} (zìxìn)}{自信 (zìxìn)}]{\zh{自信} (zìxìn) — self-confident; sure of oneself}

\label{lesson:ZH-C465-zixin}



\begin{quote}
Before the new one: say the Mandarin for optimistic, then the Mandarin for pessimistic.
\end{quote}

\subsection*{You'll want to know: \zh{自信}}
\textbf{\zh{自信}} — \emph{zìxìn} — \textquotedblleft{}self-confident; sure of oneself\textquotedblright{}.

\begin{quote}
\textbf{\zh{看到他这么自信我很高兴}} \emph{kàn dào tā zhème zìxìn wǒ hěn gāoxìng} — I'm really glad to see him so sure of himself.
\end{quote}

\begin{cousinweb}
\textbf{\zh{自信}} (\emph{zìxìn}) is built from \zh{自} 'self' and \zh{信} 'to trust', both inherited from Middle Chinese, \emph{dzijH} and \emph{sinH} (Guangyun, in Baxter's transcription), which glosses the first 'from' and the second 'faithful' and 'proof'.
\end{cousinweb}

\begin{grammarlens}[title={What you've built}]
One more word for saying how it was.
\end{grammarlens}

\subsection*{Your turn}
\begin{itemize}
\raggedright
  \item \emph{Say it:} \emph{zìxìn}
  \item \emph{Say it:} \emph{zìxìn}, once more
  \item \emph{Say it:} \emph{kàn dào tā zhème zìxìn wǒ hěn gāoxìng} — I'm really glad to see him so sure of himself.
  \item \emph{Recall:} say the Mandarin for optimistic, then the Mandarin for pessimistic, then say \emph{zìxìn} again
\end{itemize}

\begin{tcolorbox}[breakable,colback=teal!4,colframe=teal!35!black,title={Before you move on}]
What does \zh{自信} mean? (\textquotedblleft{}Self-confident; sure of oneself\textquotedblright{}.) Say it once more.
\end{tcolorbox}
